% https://tex.stackexchange.com/questions/85536/how-to-draw-nested-nodes
\documentclass[tikz,border=10pt]{standalone}
\usetikzlibrary{ext.nodes,arrows.meta,positioning,graphs}
\begin{document}
    \begin{tikzpicture}[
        ext/node picture={prefix name=false, reset graphic state},
        fixed height/.style={text depth=+0pt, text height=+.7\baselineskip},
        n/.style={
            shape=rectangle, rounded corners, draw, fixed height
        },
        b/.style={
            shape=rectangle, draw, ext/node picture, inner xsep=4mm,
            fill=blue!\inteval{50*(\tikzextnodepicture+1)}!red!35,
        },
        h/.style={above=+.333em of current bounding box, fixed height},
        node distance=3mm and 4mm,
    ]
    \NewDocumentCommand{\nbox}{O{} m d() m}{%
        \node[{
            b,#1,style/.expanded=\IfValueT{#3}{ext/node picture/baseline={(#3.base)}}
        }]{
            #4
            \node[h]{#2};
            \node[
                white, inner xsep=+0pt,
                below left=0pt and 0pt of current bounding box.north east
            ] {L\tikzextnodepicture};
        };
    }
    \nbox{Pipeline}{
        \node[n](IS){IntSrc};
        \nbox[base right=of IS, name=BPF]{BandPassFilter}(dup){
            \node[n](dup){dupl};
            \node[n, above right=of dup](HPF){HPF};
            \node[n, below right=of dup](LPF){LPF};
            \node[n, below right=of HPF](rr) {rr(1,1)};
        };
        \node[n, base right=of BPF](IP){IntPrnt};
    }
    \graph[edge=-Latex, use existing nodes]{
        IS -> dup ->[in=180] {HPF[>out= 90, <in= 90],LPF[>out=-90, <in=-90]}->[in distance=5mm, out=0] rr -> IP
    };
    \end{tikzpicture}
\end{document}